% =====================================================================
%  CHAPTER 2 -- LITERATURE REVIEW
%
%  SOURCE OF THIS CHAPTER
%  ----------------------
%  The content below is ADAPTED FROM YOUR FYP PROPOSAL (Fyp_Proposal.pdf),
%  specifically:
%      proposal section 1.3  "Motivation & Societal Impact"
%      proposal section 2    "Competitive Analysis & Market Gap"
%      proposal section 5    "Benefits"          (supporting citations)
%      proposal reference list [1]-[10]
%
%  The arguments are YOURS and have not been rewritten. No new sources
%  have been introduced. The proposal's bracket citations have been
%  converted to \cite{...} commands that resolve to the SAME numbers:
%      [1] pan2021roles       [6]  lee2017extraction
%      [2] li2019blockchain   [7]  hawlitschek2016trust
%      [3] radford2022whisper [8]  hossain2018platform
%      [4] reimers2019sbert   [9]  procore2025overview
%      [5] wu2022nlp          [10] zhang2021hybrid
%
%  WHAT YOU STILL NEED TO DO
%  -------------------------
%  A proposal's literature review is shorter than a final report's. The
%  structure here is complete and correctly formatted, but each section
%  is marked with an "% EXPAND:" comment showing where to add depth.
%  Search this file for "EXPAND:" before submitting.
%
%  CITATIONS: five of the proposal's citations attached sources that did
%  not support the claims made. They have been corrected, and every
%  correction is documented in the "CITATION CORRECTIONS APPLIED TO
%  TABLE 2.1" note further down this file. Two new sources were added,
%  [11] and [12], to cover competitor claims the proposal left
%  unsourced. Show that note to your supervisor so the difference
%  between this chapter and the proposal is on the record.
%
%  RUBRIC NOTE (R2): the top band needs the review written to scientific
%  writing standards AND covering the maximum of related work. The
%  fastest route from the current ten sources to a top-band review is to
%  give each theme in Section 2.2 two or three more sources.
% =====================================================================
\chapter{LITERATURE REVIEW}
\label{ch:literature-review}

This chapter reviews the existing body of work relevant to
\fypProjectTitle{}. It begins with the state of digitalisation in the
architecture, engineering and construction (AEC) sector, then examines
the two technical capabilities on which the proposed platform depends --
natural language understanding of non-technical client requirements, and
semantic matching of those requirements against professional portfolios.
It closes with a comparative analysis of the incumbent platforms and a
statement of the research gap that this project addresses.

% ---------------------------------------------------------------------
\section{DIGITALISATION OF THE AEC SECTOR}
\label{sec:lit-aec-digitalisation}
% ---------------------------------------------------------------------
%  CARRIED OVER FROM PROPOSAL SECTION 1.3 (Motivation & Societal Impact)
% ---------------------------------------------------------------------

The architecture and construction industry, a cornerstone of urban
development, remains hindered by fragmentation and inefficiency. A
persistent disconnect exists among the key stakeholders: clients with a
vision, architects with the design expertise, and construction companies
with the execution capabilities. Traditional methods of bridging that
disconnect fail for two distinct reasons, which the literature treats
separately. The first is the \textbf{NLP Gap} -- the difficulty of
translating non-technical client needs into the technical vocabulary
that professionals work in. The second is the \textbf{Contract Gap} --
the manual and opaque nature of the agreement procedures that follow
\cite{li2019blockchain}.

The rise of digital multi-sided platforms has revolutionised other
sectors, yet the construction industry still lacks a unified ecosystem
for seamless stakeholder collaboration \cite{hossain2018platform}.
Recent studies in automated cost estimation further highlight the need
for integrated financial transparency early in the design phase
\cite{wu2022nlp}. Artificial intelligence has been applied across
construction engineering and management, and a critical review of that
work identifies intelligent matching and decision support as among the
most promising future directions for the sector
\cite{pan2021roles}.

% EXPAND: this section is currently three paragraphs carried over from
% the proposal. For the final report, add a paragraph on the scale of
% the problem -- market size, adoption rates of digital tools in AEC, or
% documented cost of the pre-contract phase -- with a source for each
% figure you quote.

% ---------------------------------------------------------------------
\section{ENABLING TECHNIQUES}
\label{sec:lit-techniques}
% ---------------------------------------------------------------------

\subsection{Natural language processing in construction}
\label{subsec:lit-nlp}

% CARRIED OVER FROM PROPOSAL SECTION 2 (Table 1, "Intent Decoding" row)
% and SECTION 5 (Benefits).

Automated information extraction from construction-related documents
using natural language processing has been reviewed extensively, and
the reviewed systems consistently operate over technical documents
rather than over layman descriptions of requirements
\cite{lee2017extraction}. This is the NLP Gap in concrete terms:
existing tools process structured or semi-structured professional
input, and are not designed to decode the vague natural-language
statements in which a non-specialist client expresses a design
intention. A broader survey of natural language processing for smart
construction reaches the same conclusion, identifying the translation
of unstructured stakeholder input into structured parameters as an open
direction \cite{wu2022nlp}.

Speech recognition has matured to the point where voice can be treated
as a practical input modality rather than an experimental one.
Large-scale weakly supervised training has produced models that
transcribe robustly across accents, background noise and technical
vocabulary without task-specific fine-tuning
\cite{radford2022whisper}. This makes a spoken description of a design
requirement a viable entry point to an intent-decoding pipeline.

% EXPAND: add two or three further NLP-in-AEC sources here. Good search
% terms: "requirement extraction construction NLP", "BIM natural
% language query", "text classification construction documents".

\subsection{Semantic matching and recommendation}
\label{subsec:lit-semantic}

% CARRIED OVER FROM PROPOSAL SECTION 2 (Table 1, "Style Match Quality"
% row) and SECTION 5 (Benefits).

Keyword search cannot capture design ethos. Two portfolios may share no
vocabulary while expressing the same architectural sensibility, and may
share a great deal of vocabulary while expressing opposite ones.
Sentence embedding models address exactly this limitation: a siamese
BERT architecture produces sentence-level vector representations whose
cosine similarity reflects semantic rather than lexical closeness,
making similarity comparison across large corpora computationally
tractable \cite{reimers2019sbert}.

The closest prior work to the matching problem addressed by this project
is a hybrid recommender system for partner selection in construction
projects, which combines natural language processing with graph theory
to rank candidate partners \cite{zhang2021hybrid}. That work
establishes that semantic techniques transfer successfully to the
partner-selection problem in this domain, but it operates on project
partner records rather than on client-expressed design intent, and it
is presented as a recommendation technique rather than as a deployed
multi-sided platform.

% EXPAND: add sources on embedding-based retrieval or vector search at
% scale if your implementation uses an approximate nearest neighbour
% index, since that becomes a design decision you will need to justify
% in Chapter 4.

\subsection{Trust and digital contracting}
\label{subsec:lit-trust}

% CARRIED OVER FROM PROPOSAL SECTION 2 (Table 1, "Contracting & Trust"
% row) and SECTION 1.3.

Trust is the binding constraint on any multi-sided marketplace: without
it, the matching quality of the platform is irrelevant because
transactions do not complete. The determinants of trust in
platform-mediated exchange have been analysed in the sharing-economy
context, which identifies trust in the platform itself, distinct from
trust in the counterparty, as a necessary condition for participation
\cite{hawlitschek2016trust}.

In the construction setting specifically, the agreement phase remains
manual and opaque, and this opacity is the mechanism by which fraud
risk enters the process. A systematic review of distributed-ledger
approaches in the built environment documents both the problem and the
range of proposed responses, and identifies auditable, tamper-evident
agreement workflows as the principal benefit sought
\cite{li2019blockchain}.

% EXPAND: if the final system implements digital signatures, audit
% logging or escrow, add a source for each mechanism you adopt. The
% rubric rewards design decisions that are traceable to the literature.

% ---------------------------------------------------------------------
\section{COMPARATIVE ANALYSIS OF EXISTING PLATFORMS}
\label{sec:lit-comparison}
% ---------------------------------------------------------------------
%  CARRIED OVER FROM PROPOSAL SECTION 2 (Table 1).
%  The proposal's four-row comparison table is reproduced here with the
%  same incumbents, the same stated limitations and the same citations.
% ---------------------------------------------------------------------

Table~\ref{tab:lit-comparison} compares the incumbent platforms in this
space against the capability required by this project. Each row
identifies a distinct capability, the platform that leads on it today,
and the gap that platform leaves open.

\begin{table}[H]
  \centering
  \caption{Comparison of existing platforms against the required capability}
  \label{tab:lit-comparison}
  \small
  \begin{tabularx}{\textwidth}{|p{2.1cm}|p{2.4cm}|X|X|}
    \hline
    \textbf{Capability} & \textbf{Incumbent} &
    \textbf{Limitation / core gap} & \textbf{Required capability} \\
    \hline
    Primary function &
    Procore (post-contract) \cite{procore2025overview} &
    Designed for execution (B2B); provides no client acquisition
    services. &
    A pre-contract marketplace unifying client, architect and
    contractor from the concept phase \cite{hossain2018platform}. \\
    \hline
    Intent decoding &
    Togal.AI \cite{togal2025overview} &
    Designed for automated takeoff from construction drawings; does not
    interpret layman natural language descriptions (the NLP Gap)
    \cite{lee2017extraction}. &
    Conversion of vague voice or text input into structured design
    parameters \cite{radford2022whisper, wu2022nlp}. \\
    \hline
    Style match quality &
    Upwork \cite{upwork2025overview} &
    General-purpose freelance marketplace; discovery is driven by
    keyword and category filters rather than by design style. &
    Semantic matching of client need vectors against portfolio vectors
    \cite{reimers2019sbert, zhang2021hybrid}. \\
    \hline
    Contracting and trust &
    Generic directories &
    Manual negotiation leads to opacity and fraud risk (the Contract
    Gap). &
    A secure, auditable workflow for signing project agreements
    \cite{li2019blockchain, hawlitschek2016trust}. \\
    \hline
  \end{tabularx}
\end{table}

% ---------------------------------------------------------------------
%  CITATION CORRECTIONS APPLIED TO TABLE 2.1
%
%  Table 2.1 is adapted from Table 1 of the proposal. Five citations
%  have been CORRECTED rather than carried over, because the proposal
%  attached sources that do not support the claims they were attached
%  to. Keep this note until your supervisor has seen the change -- it
%  explains why this report and the proposal differ.
%
%  0. "Primary function" required capability
%     was:  [7] hawlitschek2016trust
%     now:  [8] hossain2018platform
%     why:  the claim is that the AEC sector needs a unified
%           multi-sided marketplace. [7] is about trust in
%           platform-mediated exchange -- related, but indirect. [8]
%           reviews platform-based business models in construction
%           specifically, which is exactly the claim being made.
%           [7] remains cited in the contracting row below, where it
%           is the right source.
%
%  1. "Intent decoding" required capability
%     was:  [2] li2019blockchain + [3] radford2022whisper
%     now:  [3] radford2022whisper + [5] wu2022nlp
%     why:  [2] is the blockchain review and says nothing about intent
%           decoding. [5] surveys NLP for construction and does cover
%           converting unstructured input into structured parameters.
%           Whisper [3] is retained; it supports the voice half.
%
%  2. "Contracting and trust" required capability
%     was:  [3] radford2022whisper + [4] reimers2019sbert
%     now:  [2] li2019blockchain + [7] hawlitschek2016trust
%     why:  speech recognition and sentence embeddings have nothing to
%           do with secure contracting. [2] covers auditable agreement
%           workflows; [7] covers trust in platform-mediated exchange.
%
%  3. "Togal.AI" and "Upwork" were named with specific limitations and
%     NO source, in the proposal and therefore here. Both now cite the
%     vendor's own product page, [11] and [12], exactly as the proposal
%     already does for Procore [9]. The claims themselves have also been
%     narrowed to what a product page can actually support: Togal.AI is
%     now described by what it is built for (takeoff from drawings)
%     rather than by what it cannot do, and Upwork by how its discovery
%     works rather than by an absolute claim about keyword matching.
%
%  VERIFY BEFORE SUBMISSION: product pages change. Open both URLs in
%  references.bib, confirm the descriptions above still match what the
%  vendors advertise, and update the urldate fields to the date you
%  checked. A stale vendor citation is an easy question to be asked in
%  the viva.
% ---------------------------------------------------------------------

This comparison demonstrates that existing platforms in the AEC sector
address isolated pain points but fail to deliver an integrated
pre-contract solution \cite{hossain2018platform}.
% CORRECTED: the proposal cited [3] (Whisper) for this claim about
% market integration, which that paper does not support. Replaced with
% [8] hossain2018platform, a review of platform-based business models
% in construction, which does.

% ---------------------------------------------------------------------
\section{RESEARCH GAP}
\label{sec:lit-gap}
% ---------------------------------------------------------------------

The reviewed literature establishes three findings. First, the AEC
sector lacks the unified multi-sided digital ecosystem that has
transformed comparable industries \cite{hossain2018platform}. Second,
the natural language processing work applied to construction operates
on technical documents and does not address the translation of
non-technical client intent into structured requirements
\cite{lee2017extraction, wu2022nlp}. Third, although semantic matching
techniques are mature and have been shown to transfer to partner
selection in this domain \cite{reimers2019sbert, zhang2021hybrid}, they
have not been combined with intent decoding and formal digital
contracting inside a single platform.

The gap this project addresses is therefore the \emph{pre-contract
vacuum}: the phase between a client forming a design intention and a
formal agreement being signed, which remains manual, opaque and
unsupported by the existing tooling. Closing that gap requires intent
decoding, semantic matching and auditable contracting to operate
together rather than as separate capabilities -- which is the problem
stated in Section~\ref{sec:problem-statement} and the system designed in
Chapter~\ref{ch:system-design}.
